\documentclass[twoside]{article}

\usepackage[preprint]{aistats2027}
\usepackage[round,authoryear]{natbib}
\usepackage{amsmath,amssymb,amsthm}
\usepackage{mathtools}
\usepackage{microtype}
\usepackage{url}

\renewcommand{\bibname}{References}
\renewcommand{\bibsection}{\subsubsection*{\bibname}}
\bibliographystyle{apalike}

\newtheorem{theorem}{Theorem}
\newtheorem{lemma}[theorem]{Lemma}
\newtheorem{proposition}[theorem]{Proposition}
\renewcommand{\thetheorem}{S\arabic{theorem}}
\newcommand{\R}{\mathbb R}
\newcommand{\ve}{\varepsilon}
\newcommand{\prox}{\operatorname{prox}}
\newcommand{\supp}{\operatorname{supp}}

\begin{document}
\runningauthor{}

\twocolumn[
\aistatstitle{Supplement to ``No Randomized Speedup for Weakly Convex Optimization''}
\aistatsauthor{Pahan Dewasurendra}
\aistatsaddress{Johns Hopkins University}
]

\appendix

\section{The Dead-Zone Chain}
\label{app:chain}

We use all definitions from Section~3 of the main paper. Let
$q_0(t)=0$ for $t\leq0$, $q_0(t)=3t^2-2t^3$ for $0<t<1$, and $q_0(t)=1$
for $t\geq1$. Thus $q_a(t)=q_0((t-a)/(1-a))$ with $a=1/16$.

\begin{lemma}[Smoothstep constants]
\label{lem:q}
The function $q_a$ is continuously differentiable, nondecreasing, and takes
values in $[0,1]$. Moreover,
\[
 0\leq q_a'\leq\frac85,\qquad
 \operatorname{Lip}(q_a')\leq\frac{512}{75},\qquad
 q_a(t)\leq\frac{256}{75}t_+^2.
\]
For $\psi(u)=N^{-1}\sum_jq_a(Nu_j/4)$,
\[
 \|\nabla\psi(u)\|\leq\frac25\sqrt N,
 \qquad \operatorname{Lip}(\nabla\psi)\leq\frac{32}{75}N.
\]
\end{lemma}

\begin{proof}
On $(0,1)$, $q_0'(t)=6t(1-t)$ and $q_0''(t)=6-12t$. Hence
$\max q_0'=3/2$ and $\operatorname{Lip}(q_0')=6$. Affine rescaling by
$1-a=15/16$ gives the first two constants. Since
$q_0(s)\leq3s_+^2$ and
$((t-a)/(1-a))_+\leq t_+/(1-a)$, the quadratic bound follows. The two
bounds on $\psi$ follow by the chain rule.
\end{proof}

For $W=\{w\in\R_+^M:\|w\|\leq1\}$ and $w_{M+1}=0$, the support-function
identity gives
\begin{align}
 \bar f(z)&=\max_{w\in W}F_w(z),\nonumber\\
 F_w(z)&=C(w)+\frac2N\sum_{m=1}^Mw_m\phi(z^m)\nonumber\\
 &\quad-\frac2N\sum_{m=1}^M(1-w_{m+1})\psi(z^m).
 \label{eq:variational}
\end{align}

\begin{lemma}[Analytic certificates]
\label{lem:analytic}
The function $\bar f$ is one-Lipschitz, one-weakly convex, finite-valued,
and bounded below. Its initial gap is at most $3M/N$.
\end{lemma}

\begin{proof}
Each $\phi$ is convex and $N/8$-Lipschitz. By
Lemma~\ref{lem:q}, the negative progress part in \eqref{eq:variational} has
weak-convexity modulus at most
\[
 \frac2N\frac{32N}{75}=\frac{64}{75}<1.
\]
The blocks are disjoint and the pointwise maximum preserves a common
weak-convexity modulus. The tilted part has Lipschitz constant at most
$\|w\|/4\leq1/4$. The progress part has gradient norm at most
\[
 \frac2N\sqrt{M\frac{4N}{25}}
 =\frac45\sqrt{\frac MN}\leq\frac4{5\sqrt{12}}.
\]
Their sum is below one.

The first term in the direct definition of $\bar f$ is nonnegative and every
$\psi$ lies in $[0,1]$. Thus $\bar f\geq-2M/N$. At zero, only the first phase
is positive and $\bar f(0)<2/N$. Since $M\geq2$, the initial gap is below
$2(M+1)/N\leq3M/N$.
\end{proof}

\begin{lemma}[Exact robust locality]
\label{lem:locality}
Let $\tau=1/(64N^2)$. If $|z_{m,j}|<\tau$ for every coordinate with
$\ell(m,j)>t$, then there is a neighborhood $\mathcal N$ of $z$ such that
\[
 \bar f(y)=\bar f(P_{t+1}y),\qquad y\in\mathcal N.
\]
\end{lemma}

\begin{proof}
Write
\[
 k_m=\min\{N,\max\{0,t-(m-1)L+1\}\}.
\]
Coordinates $j>k_m$ are more than level $t$ and have magnitude below
$\tau$. Since
$(N/4)\tau<1/16=a$, the function $q_a(Nu/4)$ is identically zero on a
neighborhood of each such coordinate. Hence every progress term locally
depends only on indices at most $k_m$.

Suppose $1\leq k_m<N$. For $j\geq k_m+2$, compare hidden affine facets:
\begin{align*}
 &\left(c_{k_m+1}-\frac N8z^m_{k_m+1}\right)
 -\left(c_j-\frac N8z^m_j\right)\\
 &\qquad\geq \frac1{16N}-\frac{N\tau}{4}
 =\frac{15}{256N}>0.
\end{align*}
Thus only facet $k_m+1$ can be active among hidden facets, and the strict
margin persists locally. The fan $\phi(z^m)$ depends only on indices at most
$k_m+1$.

Now suppose $k_m=0$ and $m\geq2$. Then $k_{m-1}\leq L$. Hidden progress
coordinates contribute zero, so $\psi(z^{m-1})\leq L/N\leq3/16$. Also
$\phi(z^m)\leq c_1+N\tau/8$. Therefore
\[
 \sigma_m(z)\leq-\frac1{16N}+\frac{N\tau}{8}
 =-\frac{31}{512N}<0.
\]
The score remains negative locally, so the positive-part norm does not depend
on block $m$. Its negative progress term is locally constant. If $m=1$, then
$t=-1$. The first facet is uniquely largest, by the same strict facet
comparison, and it belongs to level zero.

Fully revealed blocks already lie in $V_t$. Combining the cases shows that
every term in $\bar f$ is locally independent of coordinates beyond
$V_{t+1}$, which proves the identity.
\end{proof}

\begin{lemma}[Phase and sign certificate]
\label{lem:phase}
Let $b=q_a(3/4)=2783/3375$. If $\psi(z^M)<b$, then
$R(z)=\|[\sigma(z)]_+\|>0$. If $R(z)>0$, every
$g\in\partial\bar f(z)$ is coordinatewise nonpositive,
$\|g\|\geq1/(4\sqrt N)$, and
\[
 g^m_j<0,\ z^m_j>0\quad\Longrightarrow\quad z^m_j<\frac{13}{2N}.
 \label{eq:support-bound}
\]
\end{lemma}

\begin{proof}
Set $p_0=1$ and $p_m=\psi(z^m)$. Let $\widehat m$ be the first index with
$p_{\widehat m}<b$. Then $p_{\widehat m-1}\geq b$. Since $q_a$ is
nondecreasing, $p_{\widehat m}<q_a(3/4)$ implies that
$Nz^{\widehat m}_j/4<3/4$ for some $j$. Since $c_j\geq1$,
\[
 \phi(z^{\widehat m})>1-\frac38=\frac58.
\]
It follows that
$\sigma_{\widehat m}(z)>5/8+b-5/4=b-5/8>0$, proving $R(z)>0$.

When $R(z)>0$, the maximizing weight in \eqref{eq:variational} is unique:
\[
 w_m=\frac{[\sigma_m(z)]_+}{R(z)},\qquad \sum_mw_m^2=1.
\]
Every subgradient has blocks
\begin{align*}
 g^m&=\frac2N\left(w_mu_m
 -(1-w_{m+1})\nabla\psi(z^m)\right),\\
 &\hspace{34mm}u_m\in\partial\phi(z^m).
\end{align*}
If $w_m>0$, then $\sigma_m>0$ and
$\phi(z^m)>5/4-\psi(z^{m-1})\geq1/4$. The zero branch of $\phi$ is
inactive. Consequently, $u_m\leq0$ coordinatewise and
$\|u_m\|_1=N/8$, which gives $\|u_m\|\geq\sqrt N/8$.
The progress gradient is nonnegative. Both terms in $g^m$ are therefore
nonpositive and cannot cancel. Hence
\[
 \|g^m\|\geq\frac{w_m}{4\sqrt N},\qquad
 \|g\|^2\geq\frac1{16N}\sum_mw_m^2=\frac1{16N}.
\]

If $g^m_j<0$, either an active affine facet contributes or
$[\nabla\psi(z^m)]_j>0$. In the first case, its value is above $1/4$, so
\[
 c_j-\frac N8z^m_j>\frac14
 \quad\Longrightarrow\quad
 z^m_j<\frac8N\left(\frac{17}{16}-\frac14\right)
 =\frac{13}{2N}.
\]
In the second case, $Nz^m_j/4<1$, so $z^m_j<4/N$. This proves
\eqref{eq:support-bound}.
\end{proof}

Lemmas~\ref{lem:analytic}, \ref{lem:locality}, and \ref{lem:phase} prove
the robust-chain proposition in the main paper.

\section{Soft-Projection Calculus}
\label{app:wrapper}

Define $s_S(x)=\alpha(x)x$ with
$\alpha(x)=(1+\|x\|^2/S^2)^{-1/2}$.

\begin{lemma}[Soft projection]
\label{lem:soft}
The map $s_S$ has norm below $S$, is one-Lipschitz, and its Jacobian is
$6/S$-Lipschitz in operator norm. If $h$ is one-Lipschitz and
one-weakly convex, then $h\circ U^\top s_S$ is
$(1+6/S)$-weakly convex for every column-orthonormal $U$.
\end{lemma}

\begin{proof}
Direct differentiation gives
\[
 J_s(x)=\alpha I-\frac{\alpha^3}{S^2}xx^\top.
\]
Its eigenvalues are $\alpha$ in tangential directions and $\alpha^3$ in the
radial direction. They lie in $(0,1]$, proving one-Lipschitzness. Differentiating
$J_s$ in a unit direction and applying the triangle inequality gives
\[
 \|D J_s(x)\|_{\rm op}
 \leq\frac1S\left(3\alpha^3\frac{\|x\|}{S}
 +3\alpha^5\frac{\|x\|^3}{S^3}\right)\leq\frac6S.
\]

Set $A(x)=U^\top s_S(x)$. For $g\in\partial h(A(x))$, weak convexity gives
\begin{align*}
 h(A(y))&\geq h(A(x))+\langle g,A(y)-A(x)\rangle\\
 &\quad-\frac12\|A(y)-A(x)\|^2.
\end{align*}
The first-order Taylor remainder of $A$ has norm at most
$3\|y-x\|^2/S$. Since $\|g\|\leq1$ and $A$ is one-Lipschitz, the right side
is at least
\[
 h(A(x))+\langle J_A(x)^\top g,y-x\rangle
 -\frac{1+6/S}{2}\|y-x\|^2.
\]
This is the subgradient characterization of $(1+6/S)$-weak convexity.
\end{proof}

\begin{lemma}[Wrapper class]
\label{lem:wrapper-class}
For $B=512\sqrt M$ and $S=16B$, the function
\[
 H_U(x)=\frac14\left[\bar f(U^\top s_S(x))
 +3(\sqrt{B^2+\|x\|^2}-B)\right]
\]
is one-Lipschitz, one-weakly convex, bounded below, and has initial gap at
most $3M/(4N)$.
\end{lemma}

\begin{proof}
The first term before scaling is one-Lipschitz by composition. The radial
term is convex, nonnegative, and three-Lipschitz. Hence the factor $1/4$
gives one-Lipschitzness. Lemma~\ref{lem:soft} gives weak-convexity modulus
at most $(1+6/S)/4<1$. The radial term vanishes at zero, so adding it can
only raise the infimum. Lemma~\ref{lem:analytic} gives the gap bound.
\end{proof}

\section{No Stationarity Created by the Wrapper}
\label{app:stationarity}

\begin{lemma}[Wrapper stationarity]
\label{lem:stationarity}
Suppose
$|\langle u_{M,j},s_S(x)\rangle|<\tau$ for every
$j>r_N:=\lfloor N/2\rfloor$. Then
\[
 \|\nabla(H_U)_{1/2}(x)\|>\frac1{64\sqrt N}.
\]
\end{lemma}

\begin{proof}
Suppose otherwise. Let
$z=\prox_{\frac12 H_U}(x):=\arg\min_y\{H_U(y)+\|y-x\|^2\}$ and
$d=2(x-z)=\nabla(H_U)_{1/2}(x)$. Proximal optimality gives
$d\in\partial H_U(z)$. The derivative of $U^\top s_S$ is surjective, so the
exact subdifferential chain rule yields some
$g\in\partial\bar f(U^\top s_S(z))$ such that
\begin{equation}
 4d=J_s(z)Ug+\frac{3z}{\sqrt{B^2+\|z\|^2}}.
 \label{eq:prox-subgrad}
\end{equation}
Because $\|J_s(z)Ug\|\leq1$, if $\|z\|\geq B$ then
\[
 \|4d\|\geq\frac3{\sqrt2}-1>1.
\]
This contradicts $\|4d\|\leq1/(16\sqrt N)$. Thus $\|z\|<B$.

Set $\alpha=(1+\|z\|^2/S^2)^{-1/2}$,
$w=U^\top z$, and $y=\alpha w$. Since $s_S$ is one-Lipschitz,
\begin{align*}
 \|(y^M_{r_N+1},\ldots,y^M_N)\|
 &\leq \sqrt N\tau+\|s_S(z)-s_S(x)\|\\
 &\leq\frac1{64N^{3/2}}+\frac1{128\sqrt N}
 \leq\frac1{64\sqrt N}.
\end{align*}
Using Lemma~\ref{lem:q},
\begin{align*}
 \psi(y^M)
 &\leq\frac{r_N}{N}
 +\frac1N\frac{256}{75}\frac{N^2}{16}
 \sum_{j>r_N}(y^M_j)^2\\
 &\leq\frac12+\frac1{19200}<\frac{2783}{3375}=b.
\end{align*}
Lemma~\ref{lem:phase} now gives $g\leq0$ coordinatewise and
$\|g\|\geq1/(4\sqrt N)$.

Projecting \eqref{eq:prox-subgrad} by $U^\top$ gives
\[
 U^\top(4d)=\alpha g+cw,
 \quad
 c=\frac3{\sqrt{B^2+\|z\|^2}}
 -\frac{\alpha^3\langle w,g\rangle}{S^2}.
\]
Here $\alpha>255/256$, $\|w\|<B$, and $\|g\|\leq1$. Thus
\[
 c>\frac{2}{B}-\frac1{256B}>0,
 \qquad c<\frac{3+1/256}{B}.
\]
Let $I=\supp(g)$. By Lemma~\ref{lem:phase},
\[
 \|(w_I)_+\|\leq\frac{13\sqrt{MN}}{2\alpha N}
 =\frac{13\sqrt M}{2\alpha\sqrt N}.
\]
Since $c>0$, negative entries of $w_I$ reinforce $g_I$. Therefore
\begin{align*}
 \|4d\|&\geq\|\alpha g+cw\|\\
 &\geq\frac{\alpha}{4\sqrt N}
 -\frac{3+1/256}{512\sqrt M}
 \frac{13\sqrt M}{2\alpha\sqrt N}\\
 &>\frac1{5\sqrt N}.
\end{align*}
This gives $\|d\|>1/(20\sqrt N)>1/(64\sqrt N)$, a contradiction.
\end{proof}

Lemmas~\ref{lem:wrapper-class} and \ref{lem:stationarity} prove
the wrapper proposition in the main paper.

\section{The Conditional-Haar Reduction}
\label{app:rotation}

We state the random-rotation argument in the exact local-oracle form used here.
Let $\pi(1),\ldots,\pi(n)$ order coordinates by nondecreasing level. Let
$K_t=\{k:\ell(\pi(k))=t\}$.

\begin{lemma}[Adapted conditional-Haar bound]
\label{lem:adapted-haar}
Draw $U=(u_1,\ldots,u_n)$ uniformly from the Stiefel manifold and let $\xi$
be independent of $U$. Suppose $v_k$ is measurable with respect to
$\xi,u_1,\ldots,u_{k-1}$ and $\|v_k\|\leq S$ for every $k\in[n]$. If
\[
 d\geq n+2S^2\tau^{-2}\log(2n^2/\delta),
\]
then, with probability at least $1-\delta$,
\[
 |\langle u_j,v_k\rangle|<\tau,
 \qquad 1\leq k\leq j\leq n.
\]
\end{lemma}

\begin{proof}
Fix $k\leq j$ and condition on
$\mathcal G_{k-1}:=\sigma(\xi,u_1,\ldots,u_{k-1})$. The remaining ordered frame is Haar in the
orthogonal complement $E_k$ of these columns, whose dimension is
$d-k+1$. In particular, $u_j$ has a uniform marginal on the unit sphere in
$E_k$. The vector $v_k$ is fixed under the conditioning. If $a_k$ is its
projection onto $E_k$, then $\|a_k\|\leq S$. The standard spherical-cap
bound, also used by \citet[Appendix~B.3]{carmon2020stationary}, gives
\begin{align*}
 \Pr\bigl(|\langle u_j,v_k\rangle|\geq\tau
 \mid\mathcal G_{k-1}\bigr)
 &\leq 2\exp\left(-\frac{(d-k)\tau^2}{2S^2}\right)\\
 &\leq 2\exp\left(-\frac{(d-n)\tau^2}{2S^2}\right)\\
 &\leq\frac{\delta}{n^2}.
\end{align*}
The event is impossible when $a_k=0$. Averaging removes the conditioning.
A union bound over fewer than $n^2$ pairs $(k,j)$ proves the claim. Notice
that no conditioning on preceding good events is needed because adaptedness
holds globally.
\end{proof}

\begin{lemma}[Virtual-query rotation]
\label{lem:virtual}
Fix a randomized algorithm with random seed $\xi$. Through the hard horizon,
with probability at least $1-\delta$, its effective query $y_t=s_S(x_t)$
satisfies
\[
 |\langle u_{\pi(k)},y_t\rangle|<\tau
 \quad\text{whenever }\ell(\pi(k))\geq t,
\]
provided the dimension bound in Lemma~\ref{lem:adapted-haar} holds.
\end{lemma}

\begin{proof}
We first define a surrogate transcript on every history. At round $t$, feed
the algorithm the value and full subdifferential of the local wrapped function
obtained by replacing $\bar f(y)$ with $\bar f(P_ty)$. This response depends
only on $x_t$ and columns through level $t$. Recursively, the surrogate query
$\widetilde x_t$ and its effective point
$\widetilde y_t=s_S(\widetilde x_t)$ are measurable with respect to $\xi$ and
columns in levels below $t$.

Extend the algorithm after its output by fixed zero queries, so its surrogate
queries are defined through the hard horizon. Construct a length-$n$ virtual
sequence. If $k\in K_t$ and $t<H$, set $v_k=\widetilde y_t$. Set all later
entries to zero. Denote the ordered column $u_{\pi(k)}$ by $\widetilde u_k$.
For $k\in K_t$, the point $v_k$ depends only on columns in levels below $t$,
hence only on
$\xi,\widetilde u_1,\ldots,\widetilde u_{k-1}$. Repetition within a layer does
not use the same-layer columns already inserted in the virtual order.
Lemma~\ref{lem:adapted-haar}, applied to the reordered frame, now gives the
claimed inequalities for the surrogate points.

It remains to couple the surrogate and true transcripts. At round zero their
queries agree because both depend only on $\xi$. Suppose they agree through
the query at round $t$. The adapted-Haar event makes every coordinate of
$U^\top s_S(x_t)$ at level at least $t$ smaller than $\tau$. Applying
Lemma~\ref{lem:locality} with index $t-1$ shows that $\bar f$ agrees locally
with $\bar f\circ P_t$. Values and full subdifferentials are local quantities.
The exact subdifferential chain rule through $U^\top s_S$ therefore makes the
true oracle response equal to the surrogate response. The soft projection and
radial term are fixed and reveal no additional column. The next queries agree
by induction. Thus the surrogate inequalities hold for the true queries
through the hard horizon.
\end{proof}

Lemma~\ref{lem:virtual} proves the random-rotation lemma in the main paper.
Notice that local equality is what controls the entire subdifferential set. A
norm approximation would not suffice.

\section{Scaling and Completion}
\label{app:scaling}

\begin{lemma}[Exact scaling]
\label{lem:scaling}
If $H$ is one-Lipschitz and one-weakly convex, define
\[
 F(x)=\frac{G^2}{\rho}H\left(\frac\rho Gx\right).
\]
Then $F$ is $G$-Lipschitz and $\rho$-weakly convex, and
\[
 \nabla F_{1/(2\rho)}(x)
 =G\nabla H_{1/2}\left(\frac\rho Gx\right).
\]
\end{lemma}

\begin{proof}
The Lipschitz and weak-convexity claims follow directly by substitution. For
$y=\rho x/G$, changing variables in the proximal problem gives
\[
 \prox_{\frac1{2\rho}F}(x)=\frac G\rho\prox_{\frac12H}(y).
\]
Multiplying the displacement by $2\rho$ proves the gradient identity.
\end{proof}

\begin{proof}[Proof of Theorem 1]
Set $\kappa=1/64$ and
\[
 r=\min\{\rho\Delta/G^2,1\},\quad
 N=\left\lfloor\frac{\kappa^2G^2}{\ve^2}\right\rfloor,
 \quad M=\left\lfloor\frac{Nr}{12}\right\rfloor.
\]
It is enough to assume
$\ve^2\leq c_0\min\{G^2,\rho\Delta\}$ with
$c_0=\kappa^2/120$. Then $N\geq60$, $Nr\geq60$, and
$2\leq M\leq N/12$. Let $F_U$ be the scaled wrapper.
Lemma~\ref{lem:wrapper-class} and Lemma~\ref{lem:scaling} show that it is
$G$-Lipschitz and $\rho$-weakly convex. Its initial gap is at most
\[
 \frac{G^2}{\rho}\frac{3M}{4N}
 \leq\frac{G^2r}{16\rho}\leq\Delta.
\]

Let $H=(M-1)L+\lfloor N/2\rfloor$. Since $N\geq60$,
$L\geq N/6$, and therefore $H\geq MN/120$. Apply
Lemma~\ref{lem:virtual} with $\delta=1/4$. Given an algorithm for $F_U$,
rescaling its queries by $\rho/G$ and its oracle transcript by the inverse
factors in Lemma~\ref{lem:scaling} gives an induced algorithm for $H_U$.
On its good event, every query before $H$ satisfies the terminal projection
condition of Lemma~\ref{lem:stationarity}. Exact scaling gives
\[
 \|\nabla(F_U)_{1/(2\rho)}(x_t)\|
 >\frac{\kappa G}{\sqrt N}\geq\ve.
\]
 Index the oracle calls as $x_0,\ldots,x_{T-1}$. After $T<H$ calls, append
 the transcript-measurable output as $x_T$. It is still before level $H$, so
 the same statement applies to it.

The probability of the good event, averaged jointly over $U$ and algorithmic
randomness, is at least $3/4$. Hence some fixed frame $U$ makes the algorithm
fail with probability at least $3/4$, which rules out success probability
$1/2$. Finally,
$N\geq\kappa^2G^2/(2\ve^2)$ and
$M\geq Nr/24$, so
\[
 H\geq\frac{MN}{120}
 \geq\frac{\kappa^4}{11520}
 \frac{G^2\min\{\rho\Delta,G^2\}}{\ve^4}.
\]
Reducing the numerical constant accounts for the appended output query. This
proves the theorem.
\end{proof}

\bibliography{references}

\end{document}
